\documentclass[11pt]{article}

\usepackage[margin=1in]{geometry}
\usepackage{booktabs}
\usepackage{hyperref}
\usepackage{url}

\title{OntoLoop: A Governed, Replayable Runtime for Sovereign Language-Model Agents}

\author{
AutoLoop Contributors\\
\texttt{anomalyco/autoloop-app}\\
Draft systems paper
}

\date{\today}

\begin{document}
\maketitle

\begin{abstract}
Language-model agents are increasingly used as tool-using software operators: they
plan, call external systems, modify artifacts, learn from feedback, and continue work
across sessions. Existing agent frameworks have demonstrated strong reasoning,
tool-use, and self-improvement capabilities, but production use introduces a
different class of requirements: policy enforcement, tenant isolation, budget
control, replayability, rollback, and auditable evidence for every consequential
action. This paper presents OntoLoop, an engineering runtime for governed and
replayable language-model agents. OntoLoop decomposes agent execution into a
policy-control layer, planning layer, guarded execution layer, verification layer,
learning layer, and reporting layer. The runtime centers all provider and tool
activity around typed task envelopes, tri-state admission decisions, append-only
evidence records, deterministic replay boundaries, and release gates that block
production writes when evidence is incomplete. We describe the architecture,
contracts, implementation, and evaluation protocol of the current open-source
alpha artifact, implemented primarily in Rust with a Vue control plane and local
CLI surfaces. Rather than optimizing only for task success, OntoLoop treats agent
autonomy as a systems problem: every successful action must be attributable,
bounded, recoverable, and explainable.
\end{abstract}

\section{Introduction}

Large language models (LLMs) have moved from passive text generation toward
interactive agents that reason, call tools, inspect environments, write code, and
adapt over time. Work such as ReAct \cite{yao2022react}, Toolformer
\cite{schick2023toolformer}, Reflexion \cite{shinn2023reflexion}, Tree of Thoughts
\cite{yao2023tot}, Voyager \cite{wang2023voyager}, and SWE-agent
\cite{yang2024sweagent} shows that combining models with tools, feedback, memory,
and iterative control can greatly expand practical capability. However, capability
alone is not a runtime contract. A production agent must also answer operational
questions:

\begin{itemize}
  \item Which identity, tenant, policy, and budget authorized an action?
  \item Which tool or provider path was selected, and was it verified at the time?
  \item What evidence proves that an artifact was actually written, hashed, and
        linked to the corresponding decision?
  \item Can an operator replay a session, explain drift, and roll back a failed
        rollout?
  \item Can tests and static checks prevent direct bypass of the runtime guard?
\end{itemize}

OntoLoop is designed around these questions. We use \emph{sovereign runtime} in an
engineering sense: the local system, not the model provider, owns policy, admission,
budget, evidence, replay, and rollback boundaries. The objective is not to replace
model-level alignment methods such as Constitutional AI \cite{bai2022constitutional},
but to complement them with deterministic software contracts around the agent's
actions.

The current OntoLoop repository is an alpha engineering artifact. Its Rust workspace
contains a primary runtime crate, state and PostgreSQL adapters, an immutable
execution constitution crate, and a trust-kernel crate. The control surfaces include
a product CLI, a Node-based wrapper, and a Vue dashboard for snapshots, replay,
capability governance, and operator settings. The repository also contains layered
acceptance scripts and integration tests that encode the runtime's hard gates.

This paper makes four contributions:

\begin{enumerate}
  \item A layered architecture for governed LLM-agent execution that separates
        policy control, planning, execution, verification, learning, and reporting.
  \item A runtime contract centered on typed task envelopes, tri-state admission,
        no-bypass enforcement, evidence-linked artifact delivery, and replayable
        decision records.
  \item An implementation blueprint for integrating providers, tools, memory,
        GraphRAG-style knowledge, triggers, session replay, and operator dashboards
        without allowing direct execution paths around the guard.
  \item A release-oriented evaluation protocol that treats acceptance gates,
        rollback evidence, and replay consistency as first-class measures of agent
        readiness.
\end{enumerate}

\section{Design Goals}

OntoLoop is built for agent workloads that may touch files, tools, external
services, memory stores, and business-facing control planes. The design is guided
by five goals.

\paragraph{G1: All execution passes through one accountable gate.}
Provider calls, tool calls, MCP-style capability invocations, and artifact writes
must be mediated by the runtime kernel. The agent should not be able to call a
high-impact path through an untracked shortcut.

\paragraph{G2: Policy decisions are explicit and typed.}
Admission returns one of three states: \texttt{Allow}, \texttt{RequiresApproval},
or \texttt{Blocked}. Every decision carries a reason and evidence linkage.

\paragraph{G3: Evidence is append-only and replayable.}
Execution produces durable records that bind session identifiers, trace identifiers,
task identifiers, capability identifiers, policy versions, artifact hashes, and
decision outcomes.

\paragraph{G4: Learning is governed.}
Feedback, memory updates, skill promotion, and route optimization are treated as
runtime events subject to verification and rollback, not as invisible prompt drift.

\paragraph{G5: Release readiness is contract-driven.}
OntoLoop does not define production readiness as passing a single benchmark. It
requires layered checks: preflight, contract, domain suites, full-chain acceptance,
and release gate fusion.

\section{Architecture}

Figure~\ref{fig:flow} shows the canonical execution path as a text diagram. A user
intent or trigger enters the transport/session layer, is compiled into a structured
turn state, passes through planning and policy context, executes through guarded
capability routing, and finally emits evidence, learning updates, replay artifacts,
and observability reports.

\begin{figure}[h]
\centering
\begin{minipage}{0.95\linewidth}
\begin{verbatim}
User/Trigger
  -> Structured Transport + Session Bridge
  -> Query Turn State Machine + Context Compiler/Compactor
  -> Requirement Clarification + Constraint Freeze
  -> Policy Context + Knowledge Context
  -> Orchestrator + Capability Router
  -> Capability Admission
  -> Trust Admission Kernel
  -> Runtime Guard + Permission Mode
  -> Layered Tool Execution + Execution Fabric
  -> Evidence Tagger + Flow State Engine
  -> Trust Evidence Ledger
  -> Verifier + Audit
  -> Learning + Promotion Gate
  -> Memory / Org Knowledge Update
  -> Unified Query / Replay / Explanation
  -> Observability + Reports
\end{verbatim}
\end{minipage}
\caption{Condensed OntoLoop runtime flow.}
\label{fig:flow}
\end{figure}

\subsection{Layer Model}

OntoLoop maps high-level agent behavior onto neutral software layers:

\begin{table}[h]
\centering
\begin{tabular}{lll}
\toprule
Layer & Role & Primary modules \\
\midrule
Policy control & Policy, risk, budget, immutable constraints & \texttt{security}, \texttt{runtime}, \texttt{config} \\
Planning & Clarification, scope freeze, decomposition & \texttt{orchestration}, \texttt{agent} \\
Execution & Capability selection, tool/provider invocation & \texttt{tools}, \texttt{providers}, \texttt{runtime} \\
Resource kernel & CPU, memory, time, queue, sandbox control & \texttt{runtime}, \texttt{security} \\
Verification & Result validation, policy checks, reject loops & \texttt{runtime}, \texttt{observability} \\
Learning & Evidence, skills, causal edges, route memory & \texttt{memory}, \texttt{rag}, \texttt{hooks} \\
Reporting & Snapshot, replay, dashboard, health telemetry & \texttt{observability}, \texttt{dashboard-ui} \\
\bottomrule
\end{tabular}
\caption{OntoLoop's neutral layer decomposition.}
\label{tab:layers}
\end{table}

\subsection{Runtime Core}

The main Rust application object assembles configuration, providers, tools,
security policy, memory, RAG, sessions, transport, plugins, services, observability,
and the runtime kernel. Its CLI exposes command families for focus anchors,
knowledge export, triggers, bridges, skills, plugins, sessions, background tasks,
memory, frontend status, and system reports.

The runtime kernel owns:

\begin{itemize}
  \item resource limits and parallelism windows;
  \item MCP/tool circuit breaker thresholds and cooldowns;
  \item gate mode (\texttt{shadow}, \texttt{canary}, or \texttt{full});
  \item permission-mode defaults and policy-mode defaults;
  \item budget and quota enforcement;
  \item attestation configuration;
  \item hook runtime and layered tool execution.
\end{itemize}

\subsection{Control Plane}

The dashboard exposes local operational APIs for health, dashboard snapshots,
replay timelines, capability catalog views, capability governance, operator
settings, trigger webhooks, and server-sent event streams. The Vue control plane
renders summary strips, runtime sidebars, graph canvases, capability tables,
detail workbenches, replay timelines, and governance panels.

\section{Contracts}

OntoLoop's central mechanism is the conversion of agent behavior into explicit
contracts. The contract layer is intentionally more restrictive than a typical
agent loop because the system must preserve auditability under failure.

\subsection{Task Envelope}

Every consequential execution is represented as a task envelope:

\begin{verbatim}
TaskEnvelope {
  session_id,
  trace_id,
  task_id,
  capability_id,
  identity,
  payload,
  constraints,
  trust_plan
}
\end{verbatim}

The envelope binds an action to identity, constraints, and trace context before
the runtime can execute it. This gives downstream evidence records stable keys.

\subsection{Admission}

Admission decisions are tri-state:

\[
  d \in \{\texttt{Allow}, \texttt{RequiresApproval}, \texttt{Blocked}\}.
\]

The runtime does not treat ``not blocked'' as sufficient. A path that requires
approval must remain distinct from a path that is automatically allowed. This
distinction is important for human-in-the-loop operations, gray rollout, and
post-incident forensics.

\subsection{Production Write Contract}

For production writes, OntoLoop uses a five-part atomic contract:

\[
  W = (\texttt{state}, \texttt{event\_log}, \texttt{evidence\_ref},
       \texttt{relation}, \texttt{write\_proof}).
\]

A write is valid only when all five elements commit together or the transaction
rolls back. This prevents common failure modes where a state mutation exists
without evidence, or an event log claims an artifact that was never written.

\subsection{Artifact Delivery}

An artifact-completion claim is invalid without a write proof, hash, and evidence
reference. This is especially important for coding agents because natural-language
answers can claim completion even when no file-backed artifact exists.

\subsection{NoBypass}

NoBypass combines three enforcement layers:

\begin{enumerate}
  \item static scans for direct provider/tool/memory/MCP shortcuts;
  \item compile-time gates around contract boundaries;
  \item runtime gates around every task envelope.
\end{enumerate}

This pattern aims to catch bypass attempts before they become runtime incidents.

\section{Lifecycle}

The canonical lifecycle is:

\begin{enumerate}
  \item \texttt{INTENT\_INTAKE}: normalize input into a requirement brief.
  \item \texttt{CONSTRAINT\_FREEZE}: freeze acceptance criteria, risk profile,
        budget, and hard boundaries.
  \item \texttt{PLAN\_SYNTHESIS}: generate plans, critiques, and an executable
        route.
  \item \texttt{CAPABILITY\_ROUTING}: select only active and verified capabilities.
  \item \texttt{GUARDED\_EXECUTION}: enforce budgets, timeout, sandbox, and circuit
        breaker policy.
  \item \texttt{VERIFICATION\_GATE}: pass, iterate, or reject.
  \item \texttt{LEARNING\_COMMIT}: persist episodes, witness logs, skill updates,
        and graph deltas.
  \item \texttt{OBSERVABILITY\_EMIT}: publish snapshots, replay events, and reports.
\end{enumerate}

The state machine supports recovery edges. A verifier rejection returns the task
to planning; a circuit breaker or budget failure returns execution to planning;
policy denial moves the task to a blocked state.

\section{Memory and Knowledge}

OntoLoop includes a memory subsystem for reflexion episodes, witness logs, skill
records, causal edges, and SuperMemory-style pipelines. It also includes GraphRAG
components for entity-relation extraction, graph updates, retrieval, sharing gates,
and organization knowledge publication. The design is related to GraphRAG
\cite{edge2024graphrag} in that graph structure is used to support global context
and sensemaking, but OntoLoop focuses on operational traceability: memory updates
are linked to evidence and promotion gates.

Learning is therefore not simply prompt accumulation. A learning delta must be
derived from evidence, scored, verified, and optionally promoted. This prevents
the agent from silently converting one unverified failure into future policy.

\section{Implementation}

The current artifact is implemented as a Rust workspace with auxiliary frontend
and deployment surfaces. The inspected repository contains:

\begin{itemize}
  \item 215 Rust source files under \texttt{src/};
  \item 85 Rust integration test files under \texttt{tests/};
  \item 63 acceptance and operations scripts under \texttt{deploy/scripts/};
  \item 72 frontend source files under \texttt{dashboard-ui/src/}.
\end{itemize}

The workspace members are:

\begin{itemize}
  \item \texttt{ontoloop}: the primary runtime and CLI crate;
  \item \texttt{autoloop-state-adapter}: internal StateStore-compatible state
        adapter with in-memory and PostgreSQL bridge paths;
  \item \texttt{autoloop-postgres-adapter}: PostgreSQL-backed state and evidence
        adapter;
  \item \texttt{trustkernel}: a verified execution kernel crate housed under
        \texttt{NoumenonCore};
  \item \texttt{ontoloop-core}: immutable execution constitution contracts.
\end{itemize}

The frontend control plane is a Vue 3 and Vite application. The Node package
wrapper exposes an \texttt{ontoloop} binary that locates the Rust executable and
launches either direct message mode, subcommand passthrough, or TUI mode.

\section{Evaluation Protocol}

OntoLoop's evaluation protocol is release-oriented. A model-only benchmark can
measure whether an agent solves a task; OntoLoop additionally asks whether the
solution is authorized, bounded, attributed, replayable, and recoverable.

\subsection{Layered Acceptance}

The acceptance model contains five layers:

\begin{table}[h]
\centering
\begin{tabular}{lll}
\toprule
Layer & Scope & Blocking condition \\
\midrule
L0 & profile, config, storage readiness & failed preflight blocks all later layers \\
L1 & WAL, admission, evidence, NoBypass & any contract breach blocks release \\
L2 & sandbox, frontend CLI, storage, signals & any required domain failure blocks L3/L4 \\
L3 & integrated full-chain acceptance & \texttt{all\_passed != true} blocks release \\
L4 & fused release decision & \texttt{allow\_release=false} blocks rollout \\
\bottomrule
\end{tabular}
\caption{Layered acceptance protocol.}
\label{tab:acceptance}
\end{table}

\subsection{Suggested Metrics}

For future empirical submissions, we recommend reporting:

\begin{itemize}
  \item task success rate under guarded execution;
  \item percentage of actions with complete evidence records;
  \item replay match rate and explainable drift rate;
  \item NoBypass static/compile/runtime coverage;
  \item rollback latency after canary failure;
  \item budget reconciliation accuracy;
  \item tenant isolation violations;
  \item operator intervention rate and approval latency.
\end{itemize}

These metrics should be reported by release gate artifacts rather than prose-only
claims. The current paper intentionally does not invent benchmark numbers for the
alpha artifact.

\section{Related Work}

\paragraph{Reasoning and tool-use agents.}
ReAct interleaves reasoning traces and actions \cite{yao2022react}, while
Toolformer studies self-supervised tool-use learning \cite{schick2023toolformer}.
OntoLoop assumes such tool-using behavior is useful, but wraps it in explicit
authorization and evidence contracts.

\paragraph{Search, reflection, and lifelong learning.}
Tree of Thoughts expands inference into deliberate search over intermediate
thoughts \cite{yao2023tot}. Reflexion stores verbal feedback in episodic memory
\cite{shinn2023reflexion}. Voyager demonstrates open-ended skill acquisition in an
embodied environment \cite{wang2023voyager}. OntoLoop's learning layer is closest
in spirit to these systems, but it treats skill promotion as a governed runtime
event rather than an unconstrained memory update.

\paragraph{Software-engineering agents.}
SWE-agent shows that agent-computer interfaces matter for automated software
engineering \cite{yang2024sweagent}. OntoLoop adds contract and evidence layers
around the computer interface: artifact claims require write proofs, and execution
must pass the runtime guard.

\paragraph{Alignment and policy.}
Constitutional AI uses rule-based critique and AI feedback for model behavior
control \cite{bai2022constitutional}. OntoLoop operates at the systems layer:
policies are evaluated as runtime gates over concrete tasks, identities, tools,
budgets, and artifacts.

\paragraph{Graph-based retrieval.}
GraphRAG builds graph indexes and community summaries to improve global
query-focused summarization \cite{edge2024graphrag}. OntoLoop uses graph memory
for operational knowledge and routing evidence, with emphasis on provenance and
promotion gates.

\section{Limitations}

OntoLoop is currently an alpha engineering runtime. Several limitations remain.
First, the artifact's production hardening is uneven across provider and MCP
surfaces. Second, rich real-time replay visualization is not yet complete. Third,
the dashboard and policy administration interfaces are local-operation oriented,
not a full enterprise identity product. Fourth, quantitative benchmark results
must be generated by the release scripts and checked into an evidence package
before this paper can be treated as a results paper rather than a systems design
paper. Finally, the current repository has substantial ongoing changes; release
claims should be tied to clean commits and immutable artifacts.

\section{Conclusion}

OntoLoop frames LLM-agent autonomy as a runtime systems problem. The central idea is
simple: an agent action should not count as complete merely because a model said so
or a tool returned text. It must be authorized, bounded, executed through a guarded
path, linked to evidence, replayable, and recoverable. The current alpha artifact
implements this idea through Rust runtime contracts, typed task envelopes,
tri-state admission, evidence ledgers, no-bypass checks, session replay, governed
learning, and layered release gates. This design points toward a class of
agentic systems where autonomy and accountability are engineered together rather
than traded against each other.

\begin{thebibliography}{9}

\bibitem{yao2022react}
Shunyu Yao, Jeffrey Zhao, Dian Yu, Nan Du, Izhak Shafran, Karthik Narasimhan, and Yuan Cao.
\newblock ReAct: Synergizing Reasoning and Acting in Language Models.
\newblock \emph{arXiv preprint arXiv:2210.03629}, 2022.
\newblock \url{https://arxiv.org/abs/2210.03629}

\bibitem{schick2023toolformer}
Timo Schick, Jane Dwivedi-Yu, Roberto Dessi, Roberta Raileanu, Maria Lomeli,
Luke Zettlemoyer, Nicola Cancedda, and Thomas Scialom.
\newblock Toolformer: Language Models Can Teach Themselves to Use Tools.
\newblock \emph{arXiv preprint arXiv:2302.04761}, 2023.
\newblock \url{https://arxiv.org/abs/2302.04761}

\bibitem{shinn2023reflexion}
Noah Shinn, Federico Cassano, Edward Berman, Ashwin Gopinath, Karthik Narasimhan,
and Shunyu Yao.
\newblock Reflexion: Language Agents with Verbal Reinforcement Learning.
\newblock \emph{arXiv preprint arXiv:2303.11366}, 2023.
\newblock \url{https://arxiv.org/abs/2303.11366}

\bibitem{yao2023tot}
Shunyu Yao, Dian Yu, Jeffrey Zhao, Izhak Shafran, Thomas L. Griffiths, Yuan Cao,
and Karthik Narasimhan.
\newblock Tree of Thoughts: Deliberate Problem Solving with Large Language Models.
\newblock \emph{arXiv preprint arXiv:2305.10601}, 2023.
\newblock \url{https://arxiv.org/abs/2305.10601}

\bibitem{wang2023voyager}
Guanzhi Wang, Yuqi Xie, Yunfan Jiang, Ajay Mandlekar, Chaowei Xiao, Yuke Zhu,
Linxi Fan, and Anima Anandkumar.
\newblock Voyager: An Open-Ended Embodied Agent with Large Language Models.
\newblock \emph{arXiv preprint arXiv:2305.16291}, 2023.
\newblock \url{https://arxiv.org/abs/2305.16291}

\bibitem{bai2022constitutional}
Yuntao Bai et al.
\newblock Constitutional AI: Harmlessness from AI Feedback.
\newblock \emph{arXiv preprint arXiv:2212.08073}, 2022.
\newblock \url{https://arxiv.org/abs/2212.08073}

\bibitem{edge2024graphrag}
Darren Edge, Ha Trinh, Newman Cheng, Joshua Bradley, Alex Chao, Apurva Mody,
Steven Truitt, Dasha Metropolitansky, Robert Osazuwa Ness, and Jonathan Larson.
\newblock From Local to Global: A Graph RAG Approach to Query-Focused Summarization.
\newblock \emph{arXiv preprint arXiv:2404.16130}, 2024.
\newblock \url{https://arxiv.org/abs/2404.16130}

\bibitem{yang2024sweagent}
John Yang, Carlos E. Jimenez, Alexander Wettig, Kilian Lieret, Shunyu Yao,
Karthik Narasimhan, and Ofir Press.
\newblock SWE-agent: Agent-Computer Interfaces Enable Automated Software Engineering.
\newblock \emph{arXiv preprint arXiv:2405.15793}, 2024.
\newblock \url{https://arxiv.org/abs/2405.15793}

\end{thebibliography}

\end{document}
